\documentclass[10pt,a4paper]{amsart}
\usepackage[margin=19mm]{geometry}
\usepackage{amsmath,amssymb,mathtools}
\usepackage[scr=rsfs,cal=euler]{mathalfa}
\usepackage{xfrac}
\usepackage[unicode,hidelinks,bookmarks=false]{hyperref}
\newcommand{\ie}{i.e.\ }
\newcommand{\cf}{cf.\ }
\newcommand{\resp}{respectively\ }
\newcommand{\Subsection}{\subsection}
\providecommand{\nobreakdash}{\nobreak\hbox{-}}
\setlength{\emergencystretch}{2em}
\multlinegap0pt
\usepackage{amssymb}
\usepackage[shortlabels]{enumitem}
\newtheorem{lemma}{Lemma}
\title{Strict Convexity and Sharp Power Concavity for a Graphical $\sigma_2$-Curvature Equation}
\author{Shuning Xu}
\date{Source: arXiv:2608.20907v1 (CC BY 4.0)}
\begin{document}
\maketitle

\noindent\emph{Source and licence.} Strict Convexity and Sharp Power Concavity for a Graphical $\sigma_2$-Curvature Equation. Version arXiv:2608.20907v1 (CC BY 4.0), distributed by arXiv under the Creative Commons Attribution 4.0 International licence, http://creativecommons.org/licenses/by/4.0/, as recorded in the arXiv licence metadata for this exact version and shown on the abstract page https://arxiv.org/abs/2608.20907v1. Full licence notice: \texttt{licenses/arxiv-2608.20907-CC-BY-4.0.txt}.\par\smallskip

\noindent\textit{Editorial scope.} Counted unit: the article's lemma lem:real-ball (source0.tex lines 1482--1497). The article's complete original proof of it is delivered with it (source0.tex lines 1499--1587, one \texttt{proof} environment of 89 lines). No mathematics is written, repaired or re-derived: the statement and the proof are the article's own lines, in the article's own notation, and no retained line is substituted or re-flowed.  Retained uncounted as the source-local context the proof needs: lines 1442--1444.  Nothing was omitted from any retained span, so the package carries no omission record.  No retained line contains a citation, so the wrapper carries no bibliography and no bibliography entry.  No retained line contains a figure environment, an image-inclusion command or drawing code, so no image is delivered and the package declares no asset dependency.  Editorial formatting only, in addition: an AMS \texttt{amsart} preamble that restates the article's own declarations and macros used by the retained lines, namely the environments \texttt{lemma} on separate counters, together with the standard packages those lines need.  The retained environments print in this document as Lemma 1 in order of appearance, whereas the article prints the same items with the numbering of its own full document.  The complete native source of the article as distributed in the arXiv e-print for this exact version is bundled unchanged as \texttt{sources/208212/source0.tex}.
\begin{equation}\label{eq:real-main-1}
\mathcal E(D^2u,Du)=1.
\end{equation}

\begin{lemma}
\label{lem:real-ball}
On the unit ball $B_1$, the Dirichlet problem
\[
\mathcal E(D^2u,Du)=1
\quad\text{in }B_1,
\qquad
u=0
\quad\text{on }\partial B_1
\]
admits a smooth radial $\Gamma_2$-admissible solution $u_0$.
If
$v_0=-\sqrt{-u_0},$
then
$D^2v_0>0\,\,\text{in }B_1.$
\end{lemma}

\begin{proof}
We seek a radial solution $u_0=u_0(r)$, where $r=|x|$, and set
$\ell(r)=u_0'(r)^2$. The eigenvalues of $D^2u_0$ are $u_0''(r)$ and
$u_0'(r)/r$ with multiplicity $n-1$. Since
$Du_0=u_0'(r)e_r$, equation \eqref{eq:real-main-1} becomes
\[
(n-1)\frac{u_0'u_0''}{r}
+
\binom{n-1}{2}
\frac{u_0'^2}{r^2}(1+u_0'^2)
=
1.
\]
Equivalently,
\begin{equation}\label{eq:ball-ode}
r\ell'+(n-2)\ell(1+\ell)
=
\frac{2r^2}{n-1},
\qquad
\ell(0)=0.
\end{equation}

The smoothness condition at the center gives
\begin{equation}\label{eq:ball-center}
\ell(r)
=
\frac{2}{n(n-1)}r^2+O(r^4).
\end{equation}
Hence $\ell>0$ and $\ell'>0$ for all sufficiently small $r>0$.

We claim that in fact $\ell'(r)>0$ for every $0<r\le1$. Suppose
otherwise, and let $r_0>0$ be the first zero of $\ell'$. Differentiating
\eqref{eq:ball-ode} gives
\[
\ell'+r\ell''+(n-2)(1+2\ell)\ell'
=
\frac{4r}{n-1}.
\]
Evaluating at $r_0$ yields
\[
r_0\ell''(r_0)
=
\frac{4r_0}{n-1}>0,
\]
whereas the first-zero property of $\ell'$ requires
$\ell''(r_0)\le0$. This contradiction proves
\begin{equation}\label{eq:ell-positive}
\ell'(r)>0
\qquad\text{for }0<r\le1.
\end{equation}

Equation \eqref{eq:ball-ode} also gives
$\ell'(r)\le2r/(n-1)$ whenever $\ell\ge0$. Integrating from $0$ to $r$ and using
$\ell(0)=0$, we obtain
\[
0\le \ell(r)\le \frac{r^2}{n-1},
\]
so no finite-radius blow-up can occur. We therefore take the positive branch
$u_0'=\sqrt{\ell}$ and impose the zero boundary condition by setting
\begin{equation}\label{eq:ball-solution}
u_0(r)
=
-\int_r^1\sqrt{\ell(t)}\,dt.
\end{equation}
Then $u_0<0$ in $B_1$ and $u_0=0$ on $\partial B_1$. Moreover,
$u_0'>0$, $u_0''=\ell'/(2\sqrt{\ell})>0$, and $u_0'/r>0$.
Consequently,
\[
D^2u_0>0
\qquad\text{in }B_1.
\]
In particular, the graph of $u_0$ is $\Gamma_2$-admissible.

Finally, since
\[
D^2v_0
=
\frac{1}{2\sqrt{-u_0}}D^2u_0
+
\frac{1}{4(-u_0)^{3/2}}
Du_0\otimes Du_0,
\]
the first term is positive definite and the second is positive
semidefinite. Hence
\[
D^2v_0>0
\qquad\text{in }B_1.
\]
\end{proof}

\end{document}
